\subsection{Positive Systems, Convexity, and Separation}


% **Role:** finite-dimensional real convex geometry for positive solutions.
%
% Positive orthants, convex hulls, affine hyperplanes, separation, alternatives, and Farkas's lemma.

This subsection uses finite-dimensional real inner-product geometry to decide when linear systems have positive solutions. The guiding theme is that a missing positive solution is often detected by a separating hyperplane, hence by a dual vector certificate.

\begin{definition} \label{def:positive-vector-orders-orthants} \hlt{Positive Orthants}\\
For $\mathbf{x}=(x_1,\ldots,x_n)\in\R^n$, use the following order and orthant notation:
\begin{enumerate}[label=(\roman*)]
\setlength{\itemsep}{0pt}
\item $\mathbf{x}\geq\mathbf{0}$ iff $\mathbf{x}$ is \hlt{coordinatewise nonnegative}: $x_i\geq0$ for every $i$.
\item $\mathbf{x}>\mathbf{0}$ iff $\mathbf{x}$ is \hlt{nonzero nonnegative}: $\mathbf{x}\geq\mathbf{0}$ and $\mathbf{x}\neq\mathbf{0}$.
\item $\mathbf{x}\gg\mathbf{0}$ iff $\mathbf{x}$ is \hlt{coordinatewise positive}: $x_i>0$ for every $i$.
\item $\R_+^n\defeq\{\mathbf{x}\in\R^n:\mathbf{x}>\mathbf{0}\}$ is the \hlt{nonzero nonnegative orthant}.
\item $\R_{++}^n\defeq\{\mathbf{x}\in\R^n:\mathbf{x}\gg\mathbf{0}\}$ is the \hlt{coordinatewise positive interior}.
\end{enumerate}
\end{definition}

\begin{motivation} \label{mot:positive-solutions-probability-distributions} \hlt{Positive Null Vectors as Probabilities}\\
Let $A=[\mathbf{a}_1\ \cdots\ \mathbf{a}_n]\in\operatorname{Mat}(m,n,\R)$ and let $\mathbf{x}=(x_1,\ldots,x_n)>\mathbf{0}$ satisfy $A\mathbf{x}=\mathbf{0}$.
\begin{enumerate}[label=(\roman*)]
\setlength{\itemsep}{0pt}
\item Set $s\defeq x_1+\cdots+x_n>0$ and $\boldsymbol{\pi}\defeq s^{-1}\mathbf{x}$.
\item Then $\pi_i\geq0$ and $\sum_{i=1}^n\pi_i=1$, so $\boldsymbol{\pi}$ is a probability vector.
\item Since $A\boldsymbol{\pi}=s^{-1}A\mathbf{x}=\mathbf{0}$, equivalently $\sum\limits_{i=1}^n\pi_i\mathbf{a}_i=\mathbf{0}$.
\item If $\mathbf{x}\gg\mathbf{0}$, then $\pi_i>0$ for every $i$.
\end{enumerate}
Thus a positive null vector normalizes to probabilities under which the columns of $A$ have zero expected value.
\end{motivation}

\begin{principle} \label{pri:positive-system-row-space-view} \hlt{Positive Solutions Are Orthant Intersections}\\
For $A\in\operatorname{Mat}(m,n,\R)$, $\ker(A)=\operatorname{row}(A)^\perp$ by Corollary \ref{cor:matrix-fundamental-subspaces-orthogonal-complements}.\\
Hence $\{\mathbf{x}\in\R^n:A\mathbf{x}=\mathbf{0},\ \mathbf{x}>\mathbf{0}\} =\operatorname{row}(A)^\perp\cap\R_+^n$; and $\{\mathbf{x}\in\R^n:A\mathbf{x}=\mathbf{0},\ \mathbf{x}\gg\mathbf{0}\} =\operatorname{row}(A)^\perp\cap\R_{++}^n$.\\
Thus \hlt{positive null vectors} are precisely orthant intersections with $\operatorname{row}(A)^\perp$.
\end{principle}

\begin{definition} \label{def:convex-combinations-cones-hulls} \hlt{Convex Combinations, Cones, and Hulls}\\
Let $\mathbf{x}_1,\ldots,\mathbf{x}_k\in\R^n$.
\begin{enumerate}[label=\roman*.]
\setlength{\itemsep}{0pt}
\item A \hlt{convex combination} of these vectors is a vector $t_1\mathbf{x}_1+\cdots+t_k\mathbf{x}_k$, $t_i\geq0$, $\sum\limits_{i=1}^k t_i=1$.
\item A set $C\subseteq\R^n$ is \hlt{convex} if the line segment between any two of its points lies in $C$.
\item A set $K\subseteq\R^n$ is a \hlt{cone} if $\mathbf{x}\in K$ and $a\geq0$ imply $a\mathbf{x}\in K$.
\item \hlt{Convex hull} of finite set $S=\{\mathbf{x}_1,\ldots,\mathbf{x}_k\}$ is smallest convex subset of $\R^n$ containing $S$, denoted $\operatorname{conv}(S)$.
\end{enumerate}
Closedness and compactness use the Euclidean metric conventions from Definitions \ref{def:metric-topology-core-definitions} and \ref{def:open-cover-and-compactness}.
\end{definition}

\begin{proposition} \label{prop:convex-hull-finite-set-combinations-compact} \hlt{Finite Convex Hulls}\\
Let $S=\{\mathbf{x}_1,\ldots,\mathbf{x}_k\}\subseteq\R^n$ and $\Delta_{k-1}\defeq\left\{(t_1,\ldots,t_k)\in\R^k:t_i\geq0,\ \sum\limits_{i=1}^k t_i=1\right\}$.\\
Then $\operatorname{conv}(S) = \left\{\sum\limits_{i=1}^k t_i\mathbf{x}_i:(t_1,\ldots,t_k)\in\Delta_{k-1}\right\}$. Moreover, $\operatorname{conv}(S)$ is compact.
\begin{proof}
Let $T:\Delta_{k-1}\to\R^n$ be $T(t_1,\ldots,t_k)=\sum\limits_{i=1}^k t_i\mathbf{x}_i$, and set $C\defeq T(\Delta_{k-1})$. Taking $t_i=1$ and $t_j=0$ for $j\neq i$ gives $S\subseteq C$. If $K\subseteq\R^n$ is convex and $S\subseteq K$, then $C\subseteq K$, hence $C\subseteq\operatorname{conv}(S)$.\\
Conversely, $\Delta_{k-1}$ is convex and $T$ is affine, so $C$ is convex; since $S\subseteq C$, $\operatorname{conv}(S)\subseteq C$. Thus $\operatorname{conv}(S)=C$.\\
The simplex $\Delta_{k-1}$ is closed and bounded, hence compact by Heine--Borel, Theorem \ref{thm:heine-borel-rk}.\\
Since $T$ is continuous, $\operatorname{conv}(S)=T(\Delta_{k-1})$ is compact by Theorem \ref{thm:continuous-image-compact}.\\
\hlt{Structural Commentary}: A finite convex hull is the continuous image of a compact coefficient simplex.\\
\hlt{Key Insight}: Convexity is carried by the coefficients $t_i$, not by the ambient coordinates.
\end{proof}
\end{proposition}

\begin{proposition} \label{prop:minkowski-sum-convex-compact-closed} \hlt{Sums of Convex and Closed Sets}\\
For $X,Y\subseteq\R^n$, set $X+Y\defeq\{\mathbf{x}+\mathbf{y}:\mathbf{x}\in X,\ \mathbf{y}\in Y\}$. Then:
\begin{enumerate}[label=(\roman*)]
\setlength{\itemsep}{0pt}
\item $X,Y$ convex $\implies X+Y$ convex.
\item $X$ compact and $Y$ closed $\implies X+Y$ closed.
\end{enumerate}
\begin{proof}
\begin{enumerate}[label=(\roman*)]
\setlength{\itemsep}{0pt}
\item Assume $X,Y$ are convex. Let $\mathbf{z}_r=\mathbf{x}_r+\mathbf{y}_r\in X+Y$ for $r=0,1$, and let $0\leq t\leq1$. Then\\
$t\mathbf{z}_0+(1-t)\mathbf{z}_1 =\bigl(t\mathbf{x}_0+(1-t)\mathbf{x}_1\bigr) +\bigl(t\mathbf{y}_0+(1-t)\mathbf{y}_1\bigr)\in X+Y$. Thus $X+Y$ is convex.
\item Assume $X$ is compact and $Y$ is closed. Let $\mathbf{z}_j\in X+Y$ and $\mathbf{z}_j\to\mathbf{z}$.\\
Let $\mathbf{z}_j=\mathbf{x}_j+\mathbf{y}_j$ with $\mathbf{x}_j\in X$ and $\mathbf{y}_j\in Y$. Since $X$ is compact, a subsequence satisfies $\mathbf{x}_{j_\ell}\to\mathbf{x}\in X$.\\
Then $\mathbf{y}_{j_\ell}=\mathbf{z}_{j_\ell}-\mathbf{x}_{j_\ell}\to\mathbf{z}-\mathbf{x}\in Y$, as $Y$ is closed. Hence $\mathbf{z}=\mathbf{x}+(\mathbf{z}-\mathbf{x})\in X+Y$, so $X+Y$ is closed.
\end{enumerate}
\hlt{Structural Commentary}: Compactness lets one hold the $X$-part still enough for closedness of $Y$ to apply.\\
\hlt{Key Insight}: A limit in a sum can be split by passing to a convergent compact subsequence.
\end{proof}
\end{proposition}

\begin{definition} \label{def:affine-hyperplanes-halfspaces-separation} \hlt{Hyperplanes, Half-Spaces, and Separation}\\
Let $\mathbf{N}\in\R^n\setminus\{\mathbf{0}\}$ and $b\in\R$.
\begin{enumerate}[label=(\roman*)]
\setlength{\itemsep}{0pt}
\item The \hlt{affine hyperplane} with normal $\mathbf{N}$ and level $b$ is $\mathcal{H}(\mathbf{N},b) \defeq \{\mathbf{x}\in\R^n:\innerprod{\mathbf{N},\mathbf{x}}=b\}$.
\item The case $b=0$ gives the \hlt{linear hyperplane} $\mathcal{H}(\mathbf{N},0)=\operatorname{span}(\mathbf{N})^\perp$.
\item The \hlt{closed half-spaces} determined by $\mathcal{H}(\mathbf{N},b)$ are\\
$\mathcal{H}_+(\mathbf{N},b)\defeq\{\mathbf{x}\in\R^n:\innerprod{\mathbf{N},\mathbf{x}}\geq b\}$ and $\mathcal{H}_-(\mathbf{N},b)\defeq\{\mathbf{x}\in\R^n:\innerprod{\mathbf{N},\mathbf{x}}\leq b\}$.\\
The \hlt{open half-spaces} are obtained by replacing $\geq,\leq$ with $>,<$.
\item For $X\subseteq\R^n$ and $\star\in\{<,\leq,>,\geq\}$, use the \hlt{set-inequality convention} $\innerprod{\mathbf{N},X}\star b$ $\Longleftrightarrow$ $\innerprod{\mathbf{N},\mathbf{x}}\star b \ \forall \mathbf{x}\in X$.
\item An \hlt{orientation} of $\mathcal{H}(\mathbf{N},b)$ is $\sigma(\mathbf{N},b)\defeq(\sigma\mathbf{N},\sigma b)$ for $\sigma\in\{\pm1\}$.
\item Sets $X,Y\subseteq\R^n$ are \hlt{strictly separated} by $\mathcal{H}(\mathbf{N},b)$ if $\exists\sigma\in\{\pm1\}$ such that $\innerprod{\sigma\mathbf{N},X}<\sigma b<\innerprod{\sigma\mathbf{N},Y}$
\item Sets $X,Y\subseteq\R^n$ are \hlt{strongly separated} by $\mathcal{H}(\mathbf{N},b)$ if $\exists\sigma\in\{\pm1\}$ and $\varepsilon>0$ such that\\
$\innerprod{\sigma\mathbf{N},X}<\sigma b-\varepsilon<\sigma b+\varepsilon<\innerprod{\sigma\mathbf{N},Y}$.
\end{enumerate}
\end{definition}

\begin{remark} \label{rem:hyperplane-basic-geometry} \hlt{Hyperplane Geometry}\\
The half-spaces satisfy $\mathcal{H}_+(\mathbf{N},b)\cap\mathcal{H}_-(\mathbf{N},b)=\mathcal{H}(\mathbf{N},b)$, and two open half-spaces with $\mathcal{H}(\mathbf{N},b)$ partition $\R^n$.\\
Moreover, $\mathbf{N}\in\mathcal{H}(\mathbf{N},\norm{\mathbf{N}}^2)$, by Cauchy--Schwarz (Theorem \ref{thm:cauchy-schwarz-inequality}), $\mathbf{N}$ is closest point of this hyperplane to $\mathbf{0}$.
\end{remark}

\begin{figure}[H]
  \centering
  % Source-supported schematic: Roman, Advanced Linear Algebra, Chapter 15.
% A negative-slope subspace has a positive vector in its orthogonal complement.
\begin{tikzpicture}[
  x=0.88cm,
  y=0.86cm,
  >=Stealth,
  every node/.style={font=\small},
  axis/.style={->, gray!60, line width=0.5pt},
  subspace/.style={blue!70!black, line width=0.95pt},
  orthline/.style={teal!65!black, line width=0.65pt},
  orthray/.style={->, teal!65!black, line width=1pt},
  orthmark/.style={gray!60, line width=0.55pt},
  point/.style={circle, fill=black, inner sep=1.05pt}
]
  \path[use as bounding box] (-2.35,-1.7) rectangle (3.15,2.55);

  \coordinate (O) at (0,0);
  \coordinate (Y) at (1.12,1.68);

  \fill[green!8] (0.10,0.10) rectangle (2.6,2.08);
  \node[text=green!45!black, anchor=east] at (2.48,0.48) {$\R_{++}^2$};

  \draw[axis] (-2.05,0) -- (2.75,0) node[right] {$x_1$};
  \draw[axis] (0,-1.45) -- (0,2.22) node[above] {$x_2$};

  \draw[subspace] (-2.0,1.34) -- (2.25,-1.50)
    node[pos=0.13, above=1pt] {$S$};

  \draw[orthline] (-0.85,-1.28) -- (O);
  \draw[orthray] (O) -- (1.38,2.07);
  \node[text=teal!60!black, anchor=south east] at (0.84,1.24) {$S^\perp$};
  \node[circle, fill=teal!65!black, inner sep=1.15pt] at (Y) {};
  \node[text=teal!60!black, anchor=west] at (1.23,1.56) {$\mathbf{y}>0$};

  \draw[orthmark] (0.25,-0.17) -- (0.43,0.10) -- (0.18,0.27);
  \node[point, label=below left:{$\mathbf{0}$}] at (O) {};
\end{tikzpicture}

  \caption{In $\R^2$, a negative-slope line $S$ has $S^\perp$ meeting the strongly positive orthant.}
  \label{fig:positive-orthant-subspace-alternative}
\end{figure}

\begin{theorem} \label{thm:nearest-point-and-point-separation-convex-set} \hlt{Nearest Points and Separating a Point}\\
Let $C\subseteq\R^n$ be nonempty, closed, and convex.
\begin{enumerate}[label=(\roman*)]
\setlength{\itemsep}{0pt}
\item There exists a unique $\mathbf{N}_C\in C$ such that $\norm{\mathbf{N}_C}=\min_{\mathbf{x}\in C}\norm{\mathbf{x}}$.
\item If $\mathbf{b}\notin C$, set $D\defeq C-\mathbf{b}$ and let $\mathbf{N}\in D$ satisfy $\norm{\mathbf{N}}=\min_{\mathbf{z}\in D}\norm{\mathbf{z}}$.\\
Then $\mathbf{N}\neq\mathbf{0}$ and $\innerprod{\mathbf{N},C} \geq \norm{\mathbf{N}}^2+\innerprod{\mathbf{N},\mathbf{b}} > \innerprod{\mathbf{N},\mathbf{b}}$.\\
With $\alpha\defeq\innerprod{\mathbf{N},\mathbf{b}}+\frac{1}{2}\norm{\mathbf{N}}^2$, the hyperplane $\mathcal{H}(\mathbf{N},\alpha)$ strongly separates $\{\mathbf{b}\}$ and $C$.
\end{enumerate}
\begin{proof}
\begin{enumerate}[label=(\roman*)]
\setlength{\itemsep}{0pt}
\item Choose $\mathbf{x}_0\in C$ and $s\geq\norm{\mathbf{x}_0}$. Then $K\defeq C\cap\overline{B}_s(\mathbf{0})$ is nonempty and compact by Theorem \ref{thm:heine-borel-rk}.\\
Since $\norm{\cdot}$ is continuous, Theorem \ref{thm:extreme-value-theorem-compact} gives $\mathbf{N}_C\in K$ minimizing $\norm{\cdot}$ on $K$.\\
If $\mathbf{x}\in C$ and $\norm{\mathbf{x}}<\norm{\mathbf{N}_C}$, then $\mathbf{x}\in K$, a contradiction; hence $\mathbf{N}_C$ minimizes on $C$.\\
If $\mathbf{u}\neq\mathbf{v}$ are minimizers of norm $m$, then $\mathbf{w}\defeq(\mathbf{u}+\mathbf{v})/2\in C$ and, by the parallelogram identity, Theorem \ref{thm:parallelogram-identity-inner-product}, $\norm{\mathbf{w}}^2 =m^2-\frac{1}{4}\norm{\mathbf{u}-\mathbf{v}}^2 <m^2$, a contradiction.
\item The set $D=C-\mathbf{b}$ is nonempty, closed, and convex, so (i) gives its unique minimum-norm vector $\mathbf{N}$.\\
Since $\mathbf{b}\notin C$, $\mathbf{0}\notin D$, hence $\mathbf{N}\neq\mathbf{0}$. For $\mathbf{z}\in D$, define $\phi(t)\defeq\norm{\mathbf{N}+t(\mathbf{z}-\mathbf{N})}^2$, $0\leq t\leq1$.\\
Minimality gives $\phi'_+(0)\geq0$, hence $\innerprod{\mathbf{N},\mathbf{z}}\geq\norm{\mathbf{N}}^2$. Taking $\mathbf{z}=\mathbf{x}-\mathbf{b}$ yields $\innerprod{\mathbf{N},\mathbf{x}}\geq\norm{\mathbf{N}}^2+\innerprod{\mathbf{N},\mathbf{b}}$ for all $\mathbf{x}\in C$. For $\alpha\defeq\innerprod{\mathbf{N},\mathbf{b}}+\norm{\mathbf{N}}^2/2$ and $\varepsilon\defeq\norm{\mathbf{N}}^2/4$, then $\innerprod{\mathbf{N},\mathbf{b}} <\alpha-\varepsilon<\alpha+\varepsilon<\innerprod{\mathbf{N},C}$, so $\mathcal{H}(\mathbf{N},\alpha)$ strongly separates $\{\mathbf{b}\}$ and $C$.
\end{enumerate}
\hlt{Structural Commentary}: The nearest point in $C-\mathbf{b}$ supplies the normal vector to a separating hyperplane.\\
\hlt{Key Insight}: Minimal norm turns a metric problem into the linear inequality $\innerprod{\mathbf{N},\mathbf{z}}\geq\norm{\mathbf{N}}^2$.
\end{proof}
\end{theorem}


\begin{figure}[H]
  \centering
  % Source-supported schematic: Roman, Advanced Linear Algebra, Chapter 15.
% A point outside a closed convex set is strongly separated by a hyperplane.
\begin{tikzpicture}[
  x=0.88cm,
  y=0.86cm,
  >=Stealth,
  every node/.style={font=\small},
  convex/.style={draw=teal!65!black, fill=teal!8, line width=0.8pt},
  hyper/.style={red!70!black, line width=0.95pt},
  normal/.style={->, red!70!black, line width=0.9pt},
  guide/.style={gray!60, line width=0.5pt}
]
  \path[use as bounding box] (-0.35,-0.75) rectangle (5.85,2.85);

  \coordinate (b) at (0.55,0.25);
  \coordinate (p) at (2.35,0.85);

  \draw[convex]
    (p)
    .. controls (2.60,1.68) and (3.34,2.42) .. (4.48,2.18)
    .. controls (5.48,1.96) and (5.70,0.72) .. (4.90,0.02)
    .. controls (4.02,-0.72) and (2.76,-0.28) .. (p)
    -- cycle;
  \node[text=teal!60!black] at (4.42,1.55) {$C$};

  \draw[hyper] (0.7167,2.75) -- (1.7767,-0.43)
    node[pos=0.13, right=2pt] {$\mathcal{H}$};

  \node[circle, fill=black, inner sep=1.4pt, label=below:{$\mathbf{b}$}] at (b) {};
  \node[circle, fill=black, inner sep=1.2pt, label=above right:{$\mathbf{p}$}] at (p) {};

  \draw[normal] (b) -- (p)
    node[pos=0.52, below=5pt] {$\mathbf{N}$};

  \draw[guide] (1.57,0.59) -- (1.53,0.71) -- (1.41,0.67);
\end{tikzpicture}

  \caption{A nearest vector in $C-\mathbf{b}$ gives a normal direction and a strongly separating hyperplane.}
  \label{fig:convex-point-separating-hyperplane}
\end{figure}

\begin{theorem} \label{thm:compact-convex-set-separated-from-subspace} \hlt{Separating a Compact Convex Set from a Subspace}\\
Let $C\subseteq\R^n$ be nonempty, compact, and convex, and let $S\leq\R^n$ satisfy $C\cap S=\emptyset$.\\
Then there exists $\mathbf{N}\in S^\perp\setminus\{\mathbf{0}\}$ such that $\innerprod{\mathbf{N},C}\geq\norm{\mathbf{N}}^2>0$.\\
Consequently, $\mathcal{H}(\mathbf{N},\norm{\mathbf{N}}^2/2)$ strongly separates $S$ and $C$.
\begin{proof}
\begin{enumerate}[label=(\roman*)]
\setlength{\itemsep}{0pt}
\item Set $D\defeq S+C$. Since $S$ is closed in finite dimension, Proposition \ref{prop:minkowski-sum-convex-compact-closed} gives $D$ closed and convex.\\
Also $\mathbf{0}\notin D$, for $\mathbf{0}=\mathbf{s}+\mathbf{c}$ would imply $\mathbf{c}=-\mathbf{s}\in C\cap S$.
\item Apply Theorem \ref{thm:nearest-point-and-point-separation-convex-set} to $D$ and $\mathbf{0}$. Then $\exists\mathbf{N}\neq\mathbf{0}$ such that $\innerprod{\mathbf{N},\mathbf{s}+\mathbf{c}}\geq\norm{\mathbf{N}}^2$, $\mathbf{s}\in S,\ \mathbf{c}\in C$.\\
Fix $\mathbf{c}_0\in C$. If $a\defeq\innerprod{\mathbf{N},\mathbf{s}}\neq0$ for some $\mathbf{s}\in S$, then $\innerprod{\mathbf{N},\lambda\mathbf{s}+\mathbf{c}_0}=\lambda a+\innerprod{\mathbf{N},\mathbf{c}_0}\to-\infty$ for $\lambda$ of opposite sign and $|\lambda|\to\infty$, a contradiction. Hence $\mathbf{N}\in S^\perp$. Taking $\mathbf{s}=\mathbf{0}$ gives $\innerprod{\mathbf{N},C}\geq\norm{\mathbf{N}}^2>0$.\\
Since $\innerprod{\mathbf{N},\mathbf{s}}=0$ for all $\mathbf{s}\in S$, the choice $\varepsilon\defeq\norm{\mathbf{N}}^2/4$ gives strong separation by $\mathcal{H}(\mathbf{N},\norm{\mathbf{N}}^2/2)$.
\end{enumerate}
\hlt{Structural Commentary}: The sum $S+C$ is a closed convex set disjoint from $\mathbf{0}$, so point separation applies.\\
\hlt{Key Insight}: Since the separating inequality holds for all translates by $S$, scaling inside $S$ forces $\mathbf{N}$ to annihilate $S$.
\end{proof}
\end{theorem}

\begin{theorem} \label{thm:positive-subspace-alternatives} \hlt{Positive Orthant Alternatives for Subspaces}\\
Let $S\leq\R^n$. Then:
\begin{enumerate}[label=(\roman*)]
\setlength{\itemsep}{0pt}
\item $S\cap\R_+^n=\emptyset \Longleftrightarrow S^\perp\cap\R_{++}^n\neq\emptyset$.
\item $S\cap\R_{++}^n=\emptyset \Longleftrightarrow S^\perp\cap\R_+^n\neq\emptyset$.
\end{enumerate}
\begin{proof}
\begin{enumerate}[label=(\roman*)]
\setlength{\itemsep}{0pt}
\item[$(i)$]$(\Rightarrow)$ Let $\Delta\defeq\operatorname{conv}\{\mathbf{e}_1,\ldots,\mathbf{e}_n\}$. By Proposition \ref{prop:convex-hull-finite-set-combinations-compact}, $\Delta$ is compact and convex, and $\Delta\subseteq\R_+^n$.\\
Hence $\Delta\cap S=\emptyset$. Theorem \ref{thm:compact-convex-set-separated-from-subspace} gives $\mathbf{N}\in S^\perp\setminus\{\mathbf{0}\}$ with $\innerprod{\mathbf{N},\mathbf{e}_i}\geq\norm{\mathbf{N}}^2>0$, $1\leq i\leq n$. QED.
\item[$(i)$]$(\Leftarrow)$ If $\mathbf{x}\in S\cap\R_+^n$ and $\mathbf{N}\in S^\perp\cap\R_{++}^n$, then $\innerprod{\mathbf{N},\mathbf{x}}>0$, contradicting $\mathbf{N}\perp S$.
\item[$(ii)$]$(\Rightarrow)$ If $S=\{\mathbf{0}\}$, then $\mathbf{e}_1\in S^\perp\cap\R_+^n$. Assume $S\neq\{\mathbf{0}\}$.\\
Let $\mathbf{B}_1,\ldots,\mathbf{B}_k$ be a basis of $S$ and set $M=[\mathbf{B}_1\ \cdots\ \mathbf{B}_k]\in\operatorname{Mat}(n,k,\R)$, $\operatorname{rows}(M)=\{\mathbf{R}_1,\ldots,\mathbf{R}_n\}\subseteq\R^k$.\\
For $\mathbf{a}=(a_1,\ldots,a_n)$, $\mathbf{a}\in S^\perp$ $\Longleftrightarrow$ $\sum\limits_{i=1}^n a_i\mathbf{R}_i=\mathbf{0}$.\\
Set $K\defeq\operatorname{conv}\{\mathbf{R}_1,\ldots,\mathbf{R}_n\}$. Hence $S^\perp\cap\R_+^n\neq\emptyset$ $\Longleftrightarrow$ $\mathbf{0}\in K$.
If $\mathbf{0}\notin K$, Theorem \ref{thm:nearest-point-and-point-separation-convex-set} gives $\mathbf{u}\in\R^k$ with $\innerprod{\mathbf{u},\mathbf{R}_i}>0$ for all $i$. Then $\mathbf{v}\defeq M\mathbf{u}\in S$ has coordinates $v_i=\innerprod{\mathbf{u},\mathbf{R}_i}>0$, so $\mathbf{v}\in S\cap\R_{++}^n$, a contradiction. Thus $\mathbf{0}\in K$, and $S^\perp\cap\R_+^n\neq\emptyset$.
\end{enumerate}
\hlt{Structural Commentary}: Use the simplex and row hull as compact witnesses.\\
\hlt{Key Insight}: Missing positivity yields an orthogonal positive certificate.
\end{proof}
\end{theorem}

\begin{theorem} \label{thm:homogeneous-positive-solution-alternatives} \hlt{Homogeneous Positive Alternatives}\\
Let $A\in\operatorname{Mat}(m,n,\R)$. Exactly one alternative holds in each pair:
\begin{enumerate}[label=(\roman*)]
\setlength{\itemsep}{0pt}
\item $\ker(A)\cap\R_{++}^n\neq\emptyset$ or $\operatorname{col}(A^t)\cap\R_+^n\neq\emptyset$.
\item $\ker(A)\cap\R_+^n\neq\emptyset$ or $\operatorname{col}(A^t)\cap\R_{++}^n\neq\emptyset$.
\end{enumerate}
\begin{proof}
\begin{enumerate}[label=(\roman*)]
\setlength{\itemsep}{0pt}
\item Let $S\defeq\operatorname{row}(A)=\operatorname{col}(A^t)\subseteq\R^n$. By Principle \ref{pri:positive-system-row-space-view}, $\ker(A)=S^\perp$.\\
Moreover, $A^t\mathbf{v}\in\operatorname{col}(A^t)=S$ for every $\mathbf{v}\in\R^m$, and every vector of $S$ is of this form.\\
Hence the first pair is equivalent to $S^\perp\cap\R_{++}^n\neq\emptyset$ or $S\cap\R_+^n\neq\emptyset$.\\
Theorem \ref{thm:positive-subspace-alternatives}(i) says exactly that these two alternatives are mutually exclusive and exhaustive.
\item With the same $S$, one has $\ker(A)\cap\R_+^n=S^\perp\cap\R_+^n$, and $\{A^t\mathbf{v}:\mathbf{v}\in\R^m\}\cap\R_{++}^n=S\cap\R_{++}^n$.\\
Thus the second pair is equivalent to $S^\perp\cap\R_+^n\neq\emptyset$ or $S\cap\R_{++}^n\neq\emptyset$.\\
Theorem \ref{thm:positive-subspace-alternatives}(ii) says exactly that these two alternatives are mutually exclusive and exhaustive.
\end{enumerate}
\hlt{Structural Commentary}: The row space gives the certificate side; the nullspace gives the solution side.\\
\hlt{Key Insight}: Positive solvability is an orthogonal-complement alternative.
\end{proof}
\end{theorem}

\begin{method} \label{meth:positive-solution-certificates} \hlt{Certifying Positive Homogeneous Systems}\\
Let $A\in\operatorname{Mat}(m,n,\R)$ and $S\defeq\operatorname{row}(A)=\operatorname{col}(A^t)$.
\begin{enumerate}[label=(\roman*)]
\setlength{\itemsep}{0pt}
\item $\ker(A)\cap\R_{++}^n=\emptyset$ is certified by $\exists\,\mathbf{v}\in\R^m$ with $A^t\mathbf{v}>\mathbf{0}$.
\item $\ker(A)\cap\R_+^n=\emptyset$ is certified by $\exists\,\mathbf{v}\in\R^m$ with $A^t\mathbf{v}\gg\mathbf{0}$.
\item Existence follows from row-space avoidance,\\
$S\cap\R_+^n=\emptyset\Rightarrow\ker(A)\cap\R_{++}^n\neq\emptyset$; or $S\cap\R_{++}^n=\emptyset\Rightarrow\ker(A)\cap\R_+^n\neq\emptyset$.
\end{enumerate}
\end{method}

\begin{lemma} \label{lem:finitely-generated-nonnegative-cone-closed-convex} \hlt{Nonnegative Matrix Cone}\\
Let $A=[\mathbf{a}_1\ \cdots\ \mathbf{a}_n]\in\operatorname{Mat}(m,n,\R)$ and $C_A\defeq A(\R_{\geq0}^n)=\left\{\sum\limits_{i=1}^n t_i\mathbf{a}_i:t_i\geq0\right\}$.\\
Then $C_A$ is a closed convex cone in $\R^m$.
\begin{proof}
For $\mathbf{x}_r=A\mathbf{y}_r\in C_A$ with $\mathbf{y}_r\geq\mathbf{0}$.\\
For $0\leq t\leq1$, $a\geq0$, $t\mathbf{x}_1+(1-t)\mathbf{x}_2=A(t\mathbf{y}_1+(1-t)\mathbf{y}_2)\in C_A$, where $a\mathbf{x}_1=A(a\mathbf{y}_1)\in C_A$, $C_A$ is convex cone.\\
For closedness, choose for $\mathbf{x}\in C_A$ a representation $\mathbf{x}=\sum\limits_{i\in I}t_i\mathbf{a}_i,\qquad t_i>0$ with $I$ minimal.\\
If $\{\mathbf{a}_i:i\in I\}$ were dependent, choose $(d_i)_{i\in I}\neq0$ with $\sum\limits_{i\in I}d_i\mathbf{a}_i=\mathbf{0}$.\\
Replacing $d$ by $-d$ if necessary, some $d_i<0$; then $\lambda\defeq\min_{d_i<0}\frac{t_i}{-d_i}$ makes $t_i+\lambda d_i\geq0$ for all $i$ and zero for at least one $i$, contradicting minimality. Hence the active columns are independent.\\
Thus $C_A=\bigcup\limits_{\substack{I\subseteq\{1,\ldots,n\}\\ \{\mathbf{a}_i:i\in I\}\text{ independent}}} K_I$, where $K_I\defeq\left\{\sum\limits_{i\in I}t_i\mathbf{a}_i:t_i\geq0\right\}$.\\
For such $I$, the map $T_I:\R^I\to\operatorname{span}\{\mathbf{a}_i:i\in I\}$, where $T_I((t_i))=\sum\limits_{i\in I}t_i\mathbf{a}_i$, is a linear isomorphism, so $K_I=T_I(\R_{\geq0}^I)$ is closed in its finite-dimensional range, and the range is closed in $\R^m$.\\
Hence each $K_I$ is closed in $\R^m$, and the finite union above is closed.\\
\hlt{Structural Commentary}: Reduce $C_A$ to finitely many independent-column cones.\\
\hlt{Key Insight}: Minimal support turns generators independent, making each cone closed.
\end{proof}
\end{lemma}

\begin{theorem} \label{thm:farkas-lemma-strict-positive-system} \hlt{Farkas's Lemma}\\
Let $A\in\operatorname{Mat}(m,n,\R)$, $\mathbf{b}\in\R^m\setminus\{\mathbf{0}\}$, and $C_A\defeq\{A\mathbf{y}:\mathbf{y}\in\R^n,\ \mathbf{y}\geq\mathbf{0}\}$.\\
Exactly one of the following holds:
\begin{enumerate}[label=(\roman*)]
\setlength{\itemsep}{0pt}
\item $\mathbf{b}\in C_A$, equivalently $\exists\,\mathbf{u}>\mathbf{0}$ with $A\mathbf{u}=\mathbf{b}$.
\item $\exists\,\mathbf{v}\in\R^m$ such that $A^t\mathbf{v}\leq\mathbf{0}$ and $\innerprod{\mathbf{v},\mathbf{b}}>0$.
\end{enumerate}
\begin{proof}
\begin{enumerate}[label=(\roman*)]
\setlength{\itemsep}{0pt}
\item If $\mathbf{b}=A\mathbf{u}$ with $\mathbf{u}\geq\mathbf{0}$ and $A^t\mathbf{v}\leq\mathbf{0}$, then $\innerprod{\mathbf{v},\mathbf{b}} =\innerprod{\mathbf{v},A\mathbf{u}} =\innerprod{A^t\mathbf{v},\mathbf{u}} \leq0$.
\item Suppose $\mathbf{b}\notin C_A$. By Lemma \ref{lem:finitely-generated-nonnegative-cone-closed-convex}, $C_A$ is closed and convex. Since $\mathbf{0}\in C_A$, Theorem \ref{thm:nearest-point-and-point-separation-convex-set}, after changing sign if necessary, gives $\mathbf{v}\in\R^m$ and $\alpha\in\R$ such that $\innerprod{\mathbf{v},C_A}<\alpha<\innerprod{\mathbf{v},\mathbf{b}}$.\\
Thus $\alpha>0$ and $\innerprod{\mathbf{v},\mathbf{b}}>0$. For every $\mathbf{y}\geq\mathbf{0}$, $\innerprod{A^t\mathbf{v},\mathbf{y}} =\innerprod{\mathbf{v},A\mathbf{y}} <\alpha$.\\
If $(A^t\mathbf{v})_i>0$, taking $\mathbf{y}=\lambda\mathbf{e}_i$ and $\lambda\to\infty$ contradicts the bounded inequality. Hence $A^t\mathbf{v}\leq\mathbf{0}$.
\end{enumerate}
\hlt{Structural Commentary}: Nonrepresentability means $\mathbf{b}\notin C_A$.\\
\hlt{Key Insight}: Separation turns nonmembership in a cone into the dual inequality $A^t\mathbf{v}\leq0$.
\end{proof}
\end{theorem}

\begin{example} \label{ex:farkas-weak-inequality-necessary} \hlt{Why the Weak Inequality Matters}\\
The conclusion $A^t\mathbf{v}\leq\mathbf{0}$ cannot generally be strengthened to $A^t\mathbf{v}\ll\mathbf{0}$. Let $A=
\begin{bmatrix}
1 & 0 & 0\\
0 & 1 & -1
\end{bmatrix}$ and $\mathbf{b}=
\begin{bmatrix}
-1\\
0
\end{bmatrix}$.\\
If $\mathbf{u}\geq\mathbf{0}$, then $(A\mathbf{u})_1=u_1\geq0$, so $A\mathbf{u}\neq\mathbf{b}$.\\
The certificate $\mathbf{v}=
\begin{bmatrix}
-1\\
0
\end{bmatrix}$ satisfies $A^t\mathbf{v}=(-1,0,0)^t\leq\mathbf{0}$ and $\innerprod{\mathbf{v},\mathbf{b}}=1>0$.\\
For arbitrary $\mathbf{w}=(r,s)^t$, $A^t\mathbf{w}=(r,s,-s)^t$, so $A^t\mathbf{w}\ll\mathbf{0}$ would force $s<0<-s$, impossible.\\
\hlt{Structural Commentary}: Opposite columns force a zero in every certificate.\\
\hlt{Key Insight}: Boundary separation is sometimes the best possible separation.
\end{example}
